% Protocol frozen before the camera-calibrated, scale-unprovided TEST evaluation.
% Result-dependent claims and tables are guarded separately in results.tex.
\subsection{Unified benchmark protocol}
\paragraph{Inputs and camera calibration.}
The task provides an edited target RGB view and three RGB views of the
undeformed source object, together with their camera poses, horizontal fields
of view, and acquisition distances. We do not provide an object-specific
scale computed from either the source or the edited target mesh to a
generator. The official single-view Pixal3D control receives only the edited
target RGB and its camera calibration; its projective pipeline uses the native
unit-scale default when the scale argument is absent. Our frozen V1 and
selected multi-view methods use the same image-only scale estimator, trained
on the training split from RGB/A and camera inputs. We report the actual
inputs used by each external method; methods without a calibrated-camera
interface retain their native inference path.

\paragraph{Common population and outputs.}
We evaluate the frozen TEST508 split, report each method's success count out
of 508, and compute cross-method ranking only on the intersection of
successful cases. The new official RGB1 control and our two multi-view
methods use a 1024 cascade, whereas the previously frozen external Pixal3D
baseline uses its native 1536 cascade and estimated cameras. Each retains its
frozen stage sampling settings, which are reported. We do not force other methods to have the same
triangle count: all raw meshes are sampled with the common 20,000-point
surface protocol and rendered at 512 pixels using the same Blender/Cycles
scene, materials, and camera records. For the official single-view Pixal3D
control and our two multi-view methods, we bake a 2048-pixel texture on a
render-only mesh with a preferred budget of 250,000 faces. If native mesh
cleanup or fast simplification fails, or the baked mesh exceeds
0.01 world units of 99th-percentile nearest-neighbor discrepancy from the
untouched output, we retry at 500,000 and then one million faces;
exceptionally complex meshes may then use two or five million faces under the
same acceptance rule. We record the
actual budget per example. The untouched 1024
output is always used for geometric metrics; this bake does not alter the
generation cascade. At budgets of one million faces or higher, we pre-simplify with
\texttt{fast-simplification} 0.2.0 and preserve that render mesh's topology
during PBR baking; the same rule is used for the official and multi-view
Pixal3D controls. Every baked mesh must pass the same 0.01 world-unit check.
Other methods retain their native textured exports under
the common renderer's one-million-face render-only cap. The official
single-view Pixal3D and
the edited-target multi-view pipelines are labeled separately because their
image inputs differ.

\paragraph{Metrics and optional source-mesh registration.}
Source preservation uses LPIPS, DreamSim, and CLIP cosine on cameras 9 and 10.
Global naturalness uses FID and KID on the same common-success image pool.
Target alignment uses PSNR, SSIM, and foreground mask IoU at camera 8, plus
symmetric unsquared Chamfer distance and F-score at threshold 0.02 on the
surface samples. The surface frame is normalized by the input source object's
center and extent, with no per-output alignment in the raw results. A second,
paired condition applies a similarity transform fitted by trimmed symmetric ICP
to the available undeformed source mesh \emph{after} generation. Registration
never uses the edited target mesh, depth, mask, or rendering, and the aligned
textured outputs are rerendered before image metrics are recomputed. The
source-aligned condition forms the main comparison; paired raw results show
the effect of registration separately for target-camera PSNR, SSIM, and mask
IoU on the exact same success subset. We mark the extra source-mesh input
explicitly. Confidence intervals resample source clusters, retaining all
edited examples from a sampled source together. Neither TEST labels nor TEST
metrics choose a scale estimator, model checkpoint, or ICP hyperparameter.

% Development evidence only; do not treat the small native-render PBR panel as
% the final common Cycles evaluation or promote old M1-scale tables here.
% Evidence: docs/superpowers/reports/2026-10-28-scale-ungiven-VAL31-PBR8-paired.json
\paragraph{Scale-unprovided development gate.}
We first hold the image-only scale estimator and the downstream Shape512,
Shape1024, and vanilla PBR weights fixed while changing only the SS stage.
On 31 source-disjoint validation edits, the V1, Retarget, and Selected
variants obtain CD $0.012610/0.012435/0.012367$ and F@0.02
$0.848527/0.850373/0.853088$, respectively. The paired F-score changes
relative to V1 are $+0.001846$ and $+0.004561$; source-cluster bootstrap
95\% intervals are $[-0.007848,+0.011855]$ and
$[-0.005920,+0.016321]$. The intervals include zero, so this panel does not
establish a geometric gain from either module. Its SS occupancy macro-IoU
is $0.443674/0.444928/0.442446$; stage occupancy and decoded mesh quality
are reported separately.

For appearance isolation, we use the same eight image-scale V1 meshes,
initial PBR noise, and native render references for vanilla and Top-1 R PBR.
Their source-camera LPIPS is $0.070836/0.070935$, and target-camera PSNR is
$24.5466/24.5072$ dB. The paired Top-1-minus-vanilla LPIPS and PSNR
intervals are $[-0.000208,+0.000480]$ and $[-0.1018,+0.0174]$ dB.
This small native-render panel does not establish appearance improvement;
the final Blender/Cycles comparison and source-M0 postprocessing remain
separate evaluations.
